\documentclass[10pt,a4paper]{amsart}
\usepackage[margin=19mm]{geometry}
\usepackage{amsmath,amssymb,mathtools}
\usepackage[scr=rsfs,cal=euler]{mathalfa}
\usepackage{xfrac}
\usepackage[unicode,hidelinks,bookmarks=false]{hyperref}
\newcommand{\ie}{i.e.\ }
\newcommand{\cf}{cf.\ }
\newcommand{\resp}{respectively\ }
\newcommand{\Subsection}{\subsection}
\providecommand{\nobreakdash}{\nobreak\hbox{-}}
\setlength{\emergencystretch}{2em}
\multlinegap0pt
\usepackage{bm}
\usepackage{bbm}
\usepackage{dsfont}
\usepackage{xspace}
\usepackage{cancel}
\usepackage{mathtools}
\usepackage{amsthm}
\makeatletter
\@ifundefined{theorem}{\newtheorem{theorem}{Theorem}}{}
\@ifundefined{lemma}{\newtheorem{lemma}[theorem]{Lemma}}{}
\makeatother
\newcommand{\LL}{\mathcal{L}}
\newcommand{\R}{\mathbb{R}}
\title[On the boundedness of some real line arrangements of type at]{On the boundedness of some real line arrangements of type at most one}
\author{Marek Janasz}
\date{Source: arXiv:2602.05409v1 (CC BY 4.0)}
\begin{document}
\maketitle

\noindent\emph{Source and licence.} On the boundedness of some real line arrangements of type at most one. Version arXiv:2602.05409v1 (CC BY 4.0), distributed under the Creative Commons Attribution 4.0 International licence, as recorded in the arXiv licence metadata for this exact version and shown on the abstract page https://arxiv.org/abs/2602.05409v1. Full rights notice: licenses/arxiv-2602.05409-CC-BY-4.0.txt.\par\smallskip

\noindent\textit{Editorial scope.} Counted unit: the Lemma whose source label is \texttt{thm:tk\_global\_bound} in the article (source0.tex lines 224--250, source label \texttt{thm:tk\_global\_bound}), delivered together with the article's own complete original proof of it (source0.tex lines 252--338, one \texttt{proof} environment of 87 lines). No mathematics is written, repaired or re-derived: statement and proof are the article's own text line for line. Required context retained uncounted: none. Every definition, display and label of those lines is retained unchanged. Omitted from this package and not redistributed: 1--223, 251--251, 339--479. Nothing inside a retained excerpt is omitted or altered: every retained body is delivered exactly as the article's lines are written, so no omission receipt of any kind accompanies this package. Citations: the retained excerpts carry no citation command at all, so no bibliography entry of the article is transcribed and no reference is redistributed. Reference adaptation: none was needed and none was made; every reference inside the delivered unit resolves inside it, so no unresolved reference or citation is produced. Printed slips kept literal: no printed word, symbol, hypothesis or number of the counted statement or of its proof was altered. Layout and class: the article's own class is \texttt{article} with packages including geometry, setspace, fontenc, lmodern, microtype, amsmath, amssymb, amsthm, mathtools, mathrsfs, pgf, tikz, natbib, hyperref; the wrapper is the lane's pinned \texttt{amsart} editorial wrapper at 10pt on A4, which restates the theorem-like environments the retained text needs and adds the article's own macro definitions that the retained text needs (\texttt{\textbackslash LL}, \texttt{\textbackslash R}), with extra wrapper packages (bm, bbm, dsfont, xspace, cancel, mathtools). The delivered numbering is the unit's own. Assets: no retained span contains a figure environment, a graphics-inclusion command, a PDF-page inclusion, a table or a TikZ picture; no image file is copied into this package, the package declares no asset dependency and no image file is redistributed with it. Licence: version arXiv:2602.05409v1 of this article is distributed under the Creative Commons Attribution 4.0 International licence, as recorded in the arXiv licence metadata for this exact version and shown on the abstract page https://arxiv.org/abs/2602.05409v1. A full rights notice is delivered at licenses/arxiv-2602.05409-CC-BY-4.0.txt. Characters: every retained excerpt is pure ASCII, so the delivered engine, XeLaTeX with its default Latin Modern text font, reports no missing character.
\begin{lemma}\label{thm:tk_global_bound}
Let $\mathcal \LL$ be an arrangement of $d$ lines in $\mathbb{P}^2_\R$.
Fix an integer $k\ge 4$ and assume that all intersection points have multiplicity at most $k$,
i.e.\ $t_i(\mathcal L)=0$ for all $i\ge k+1$. Then the number of $k$--fold
intersection points $t_{k}$ satisfies the following bounds.

\smallskip
\noindent\textbf{(1) Case $d\ge  k+3$.} One has
\[
t_k \;\le\; \frac{d(d-1)-16}{k^2+3k-15}.
\]
\smallskip
\noindent\textbf{(2) General case.} One has
\[
t_k \;\le\; \frac{d}{k}\left\lfloor\frac{d-1}{k-1}\right\rfloor .
\]
\smallskip
In particular, for every $d\ge 2$ we have
\[
t_k \;\le\;
\max\!\left\{
\frac{d}{k}\left\lfloor\frac{d-1}{k-1}\right\rfloor,\;
\left\lfloor\frac{d(d-1)-16}{k^2+3k-15}\right\rfloor
\right\},
\]
where the second term makes sense only when $d\ge k+3$.
\end{lemma}

\begin{proof}
Counting unordered pairs of distinct lines in two ways yields
\begin{equation}\label{eq:pairs_general_k}
\binom{d}{2} \;=\; \sum_{r\ge 2}\binom{r}{2}\,t_r.
\end{equation}
Under the assumption $t_r=0$ for $r\ge k+1$, identity \eqref{eq:pairs_general_k} becomes
\begin{equation}\label{eq:pairs_trunc_general_k}
\binom{d}{2} \;=\; t_2 + \sum_{r=3}^{k}\binom{r}{2}\,t_r,
\end{equation}
hence
\begin{equation}\label{eq:t2_general_k}
t_2 \;=\; \binom{d}{2} - \sum_{r=3}^{k}\binom{r}{2}\,t_r.
\end{equation}

\medskip
\paragraph{The case $d\ge k+3$.}
Then the maximal multiplicity satisfies $m(\mathcal L)\le k\le d-3$, so we may apply Shnurnikov's inequality
\begin{equation}\label{eq:HL_general_k}
t_2 + \frac{3}{2}t_3 \;\ge\; 8 + \sum_{r\ge 4}\frac{4r-15}{2}\,t_r.
\end{equation}
Since $t_r=0$ for $r\ge k+1$, \eqref{eq:HL_general_k} reduces to
\begin{equation}\label{eq:HL_trunc_general_k}
t_2 + \frac{3}{2}t_3 \;\ge\; 8 + \sum_{r=4}^{k}\frac{4r-15}{2}\,t_r.
\end{equation}
Substituting \eqref{eq:t2_general_k} into \eqref{eq:HL_trunc_general_k} gives
\[
\binom{d}{2}
-
\sum_{r=3}^{k}\binom{r}{2}\,t_r
+
\frac{3}{2}t_3
\;\ge\;
8 + \sum_{r=4}^{k}\frac{4r-15}{2}\,t_r,
\]
hence, after rearranging,
\begin{equation}\label{eq:rearranged_general_k}
\binom{d}{2} - 8
\;\ge\;
\frac{3}{2}t_3
+
\sum_{r=4}^{k}
\left(
\binom{r}{2} + \frac{4r-15}{2}
\right)t_r.
\end{equation}
All terms on the right-hand side are nonnegative. In particular,
\[
\binom{d}{2} - 8
\;\ge\;
\left(
\binom{k}{2} + \frac{4k-15}{2}
\right)t_k.
\]
A direct computation yields
\[
\binom{k}{2} + \frac{4k-15}{2}
=
\frac{k(k-1)}{2} + \frac{4k-15}{2}
=
\frac{k^2+3k-15}{2}.
\]
Therefore,
\[
t_k
\;\le\;
\frac{2\bigl(\binom{d}{2}-8\bigr)}{k^2+3k-15}
=
\frac{d(d-1)-16}{k^2+3k-15},
\]
which proves (1).

\medskip
\paragraph{General bound.}
Fix a line arrangement $\LL$. For every line $\ell$ in $\LL$, the following obvious count holds
\[(d-1) = \sum_{p \in {\rm Sing}(\mathcal{L}) \cap \ell} (r-1)t_{r,\ell},\]
where $t_{r,\ell}$ denotes the number of $r$-fold points located on the line $\ell$.
The above formulate tells us that each $k$--fold point from ${\rm Sing}(\mathcal{L}) \cap \ell$ costs $k-1$ intersections (since $k-1$ other lines meet $\LL$ at that point). Hence the number of $k$--fold points in ${\rm Sing}(\mathcal{L}) \cap \ell$ is at most $\lfloor (d-1)/(k-1)\rfloor$.
Counting incidences between lines and $k$--fold points gives us
\[
k\,t_k \;\le\; d\left\lfloor\frac{d-1}{k-1}\right\rfloor,
\]
and thus
\[
t_k \;\le\; \frac{d}{k}\left\lfloor\frac{d-1}{k-1}\right\rfloor,
\]
which proves (2). The final unified estimate follows immediately.
\end{proof}

\end{document}
